\clearpage
\pdfbookmark[1]{Fourier. 25 February 1833.}{anzeige_fourier_128_130}
\SourceLeaf{128}
\begin{center}
\rule{0.55\linewidth}{0.4pt}\par Göttingische gelehrte Anzeigen.\quad 25 February 1833.\par \rule{0.55\linewidth}{0.4pt}
\end{center}
\textit{Analysis of determinate equations, by} \textsc{Fourier}.\quad\textit{First part.}\quad Paris 1831. \textit{Published by Firmin Didot Brothers.}\quad (XXIV and 258 pages in quarto.)

According to the author's plan, this work was to consist of two divisions. The first was to contain a general survey of the whole and the first two sections, while the second division was to contain the remaining five sections. Death tore the author away from the sciences just as printing had begun. The whole first division, however, which we receive here through the efforts of Mr \textsc{Navier}, was found almost fully worked out, whereas the rest was still far from completion. This is all the more regrettable because the general survey shows that several interesting investigations were intended for it. Passing over what we might have expected had the author been granted a longer life, we must here confine our account to what we have actually received.

The purpose of the work is to give the numerical solution of determinate algebraic equations an entirely reliable and as simple a course as possible: in particular, to ascertain the number of real and imaginary roots with complete certainty, to enclose the former individually within fixed bounds, and to arrange entirely methodically the successive approximation to their values, which can be carried to any desired accuracy. The principal foundation of the whole is a theorem of the author's own,
\SourceLeaf{129}
which is to be regarded as a successful generalization of \textsc{Descartes}' theorem (usually named after \textsc{Harriot}) and consists essentially in the following. Let $X$ be an ordered polynomial in $x$ (with all its numerical coefficients real). Further, let $a,a'$ be any two fixed, unequal real quantities, with $a'-a$ positive. Under the substitutions $x=y+a,\ x=y+a'$, let $X$ become the likewise ordered functions $Y,Y'$, whose coefficients contain $g,g'$ sign continuations, respectively. The exceptional case in which one or more coefficients vanish is set aside here for brevity. Under these assumptions, the equation $X=0$ cannot, within the bounds $x=a,\ x=a'$, contain more than $g'-g$ real roots; or, more precisely, if the number of real roots between those bounds is $=\lambda$, then $g'-g-\lambda$ is always either $=0$ or a positive even number. Consequently, if the difference $g'-g$, which can never be negative, is $=0$, there is no real root at all between those bounds; if $g'-g=1$, the bounds contain one real root and no more; finally, if $g'-g=2$, it remains uncertain for the time being whether two real roots lie between the bounds or none at all. Instead of two bounds $a,a'$, a larger number can be treated in a very similar way and chosen so that, first, the smallest corresponds to a function $Y$ with sign changes only, and the largest to one with sign continuations only. Thus, if the corresponding numbers of sign continuations are denoted in order by $g,g',g'',g'''$, and so on, the first of these numbers will be $=0$ and the last equal to the degree of the equation. Second, all the individual differences $g'-g,\ g''-g',\ g'''-g''$, etc., are to be either $=1$ or $=2$ only. All the real roots are thereby enclosed between bounds so that each individual interval must contain one root or may contain two. How, in the latter case, an entirely methodical decision is reached, if necessary by further narrowing the bounds, as to whether the two real roots are actually present or absent cannot be explained further here because of the limited space. We note only that whenever this last case occurs, there is always an intermediate value within such bounds for which a coefficient before the last in the function $Y$ vanishes, while the preceding and following coefficients have the same sign. \textsc{Fourier} calls such positions critical. Each such critical
\SourceLeaf{130}
position therefore entails the absence of two real roots. When, however, \textsc{Fourier} also expresses himself by saying that two such absent real roots become \textit{imaginary} each time, we cannot entirely approve this expression, since it could easily give rise to a misinterpretation. It is indeed true that, \textit{counted together}, the equation $X=0$ contains exactly as many pairs of imaginary roots as there are such absences or critical positions. But the values of all the imaginary roots are in themselves just as determinate as the real ones. That expression can therefore easily be interpreted as saying that each \textit{particular} gap has a \textit{particular} pair of imaginary roots belonging to it. Not only has \textsc{Fourier} \textit{not} demonstrated this; it must remain in doubt until deeper investigations have shed light on this interesting point. This is not meant to say that \textsc{Fourier} himself understood the expression in that way. We would rather assume the opposite and almost suspect that he remained uncertain whether such a definite connection exists or not, and deliberately avoided an open declaration about the potentially misleading expression. In general, \textsc{Fourier} has not entered into the nature and calculation of imaginary roots at all in this work, and a wide field therefore remains to be worked.\par
For so skilled a mathematician as \textsc{Fourier}, once in possession of that beautiful theorem---and, according to the notes communicated by Mr \textsc{Navier}, he had been in possession of it for a very long time---it could not have been difficult to base the organization of the technique for the numerical solution of equations on it. This development is given with great completeness and detail. More practised readers might perhaps prefer a somewhat more condensed presentation and the removal of some repetitions; the less practised will welcome the many well-chosen and fully treated examples. In any case, this work secures \textsc{Fourier}'s name an honourable place in this branch of mathematics too, such as he has long held in others.\par
\par\noindent\hfill\rule{0.30\linewidth}{0.4pt}\hfill\mbox{}\par
